% 数学ノート集・共通テンプレート v2（積分論2026年度第1部の版面を参考）
% A5判。章内で定義・定理・例・注意等の番号を共通化。
% コンパイル： platex -interaction=nonstopmode math_note_template_v2.tex (2回)
%              dvipdfmx math_note_template_v2.dvi
\documentclass[a5paper,10pt,onecolumn]{jsarticle}
\usepackage[dvipdfmx,top=22mm,bottom=20mm,inner=19mm,outer=18mm,headheight=13pt,headsep=6mm]{geometry}
\usepackage{amsmath,amssymb,amsthm,mathtools}
\usepackage{enumitem,fancyhdr,titlesec}
\usepackage[dvipdfmx,bookmarks=true,colorlinks=false]{hyperref}
\usepackage{pxjahyper}

% ----- 差し替える文書情報 -----
\newcommand{\NoteTitle}{滑らかな多様体・接空間・微分}
\newcommand{\NoteSeries}{数学ノート集}

% ----- 文字・組版 -----
% 参考PDFの質素な明朝＋Computer Modernの印象を引き継ぎながら，
% A5判で窮屈にならないよう本文の行間を少し確保する．
\linespread{1.13}
\setlength{\parindent}{1zw}
\setlength{\parskip}{0pt}
\setlength{\abovedisplayskip}{6pt plus 2pt minus 2pt}
\setlength{\belowdisplayskip}{6pt plus 2pt minus 2pt}
\setlength{\abovedisplayshortskip}{3pt plus 2pt minus 1pt}
\setlength{\belowdisplayshortskip}{4pt plus 2pt minus 1pt}
\setlist[enumerate]{topsep=3pt,itemsep=1pt,parsep=0pt,leftmargin=2.4zw}
\titleformat{\section}{\normalfont\normalsize\bfseries}{\thesection}{0.8zw}{}
\titlespacing*{\section}{0pt}{2.0ex plus .4ex}{.9ex}
\titleformat{\subsection}{\normalfont\normalsize\bfseries}{\thesubsection}{0.7zw}{}
\titlespacing*{\subsection}{0pt}{1.5ex plus .3ex}{.6ex}
\raggedbottom

% ----- 柱とページ番号：ページ上部、罫線なし -----
\fancypagestyle{mathnote}{%
  \fancyhf{}%
  \fancyhead[L]{\footnotesize \NoteSeries\ （\NoteTitle）}%
  \fancyhead[R]{\footnotesize \thepage}%
  \renewcommand{\headrulewidth}{0pt}%
  \renewcommand{\footrulewidth}{0pt}%
}
\pagestyle{mathnote}

% ----- 節内通し番号の定義・定理群 -----
\newtheoremstyle{mathnotetheorem}
 {4.5pt}{4.5pt}{\normalfont}{}%
 {\bfseries}{}{0.65zw}%
 {\thmname{#1}\thmnumber{ #2}\thmnote{（#3）}}
\theoremstyle{mathnotetheorem}
\newtheorem{theorem}{定理}[section]
\newtheorem{definition}[theorem]{定義}
\newtheorem{lemma}[theorem]{補題}
\newtheorem{proposition}[theorem]{命題}
\newtheorem{corollary}[theorem]{系}
\newtheorem{example}[theorem]{例}
\newtheorem{remark}[theorem]{注意}
\renewcommand{\proofname}{証明}
\renewcommand{\qedsymbol}{$\square$}
\makeatletter
\renewenvironment{proof}[1][\proofname]{%
  \par\pushQED{\qed}\normalfont\topsep4\p@\@plus2\p@\relax
  \trivlist\item[\hskip\labelsep\bfseries #1]\ignorespaces
}{\popQED\endtrivlist\@endpefalse}
\makeatother

% ----- 使うノートで適宜拡張する記号 -----
\newcommand{\R}{\mathbb{R}}
\newcommand{\norm}[1]{\lVert #1\rVert}
\newcommand{\abs}[1]{\lvert #1\rvert}

\begin{document}
\section*{読了後にできること}
本稿を読むための前提は，滑らかな多様体，接空間，正則値定理．
本稿の具体的な到達課題は，方程式で与えられた部分多様体の接空間を計算する．

\section{滑らかな多様体の定義}

本稿では $n$ 次元の多様体はHausdorffかつ第二可算であり，
滑らかさは $C^\infty$ 級を意味する．
Euclid空間の逆写像定理・陰関数定理・階数定理を前提とする．

\begin{definition}
位相空間 $M$ の各点 $p$ が
$\mathbb R^n$ の開集合と同相な開近傍 $U$ を持つとき，
$M$ を $n$ 次元\emph{位相多様体}という．
同相写像 $\varphi:U\to\varphi(U)\subset\mathbb R^n$ を
\emph{座標写像}，対 $(U,\varphi)$ を座標近傍という．
\end{definition}

\begin{definition}
座標近傍 $(U,\varphi),(V,\psi)$ が滑らかに両立するとは，
$U\cap V\ne\varnothing$ のとき
\[
 \psi\circ\varphi^{-1}:
 \varphi(U\cap V)\longrightarrow\psi(U\cap V)
\]
およびその逆写像が滑らかであることをいう．
互いに両立する座標近傍が全体を覆う集合を
\emph{滑らかなAtlas}という．
最大の滑らかなAtlasを滑らかな構造という．
\end{definition}

\begin{proposition}[最大Atlasの存在と一意性]
\label{prop:maxatlas}
任意の滑らかなAtlas $\mathcal A$ に対して，
$\mathcal A$ の全ての座標近傍と滑らかに両立する
座標近傍全体は，$\mathcal A$ を含む唯一の
最大Atlasをなす．
\end{proposition}
\begin{proof}
$\mathcal A$ と両立する座標近傍全体を
$\overline{\mathcal A}$ とする．
任意の二つ $(U,\varphi),(V,\psi)$ を取る．
$U\cap V$ の点 $p$ に対して
$\mathcal A$ の座標近傍 $(W,\chi)$ で $p\in W$ となるものが
存在する．
\[
 \psi\circ\varphi^{-1}
 =(\psi\circ\chi^{-1})\circ(\chi\circ\varphi^{-1})
\]
は $p$ の座標近傍で滑らかである．
その逆も同様なので二つの座標は両立する．
従って $\overline{\mathcal A}$ はAtlasとなる．
もしより大きい両立するAtlasがあれば，
その全ての元は $\mathcal A$ と両立するため
$\overline{\mathcal A}$ に既に属する．
ゆえに最大性と一意性を得る．
\end{proof}

\section{滑らかな写像}

\begin{definition}
滑らかな多様体 $M,N$ と連続写像 $F:M\to N$ に対し，
任意の $p\in M$ で座標近傍 $(U,\varphi)$，
$(V,\psi)$ を選び，
$F(U)\subset V$ となるように縮めた後
\[
 \psi\circ F\circ\varphi^{-1}
\]
がEuclid空間の開集合上の滑らかな写像になるとき，
$F$ を滑らかという．
\end{definition}

\begin{proposition}[滑らかさの座標独立性]
\label{prop:map-independence}
滑らかな写像であることは，
その判定に用いる座標近傍の選択によらない．
また滑らかな写像の合成は滑らかである．
\end{proposition}
\begin{proof}
別の座標近傍を $(U',\varphi')$，
$(V',\psi')$ とすると重なりで
\[
 \psi'\circ F\circ(\varphi')^{-1}
 =(\psi'\circ\psi^{-1})
 \circ(\psi\circ F\circ\varphi^{-1})
 \circ(\varphi\circ(\varphi')^{-1})
\]
と表される．
両端は滑らかな座標変換であり，中間は仮定により滑らかだから，
合成も滑らかである．
同様に $F:M\to N$ と $G:N\to P$ の座標表示は
Euclid空間の滑らかな写像の合成となるので滑らかである．
\end{proof}

\section{曲線による接空間}

\begin{definition}
$p\in M$ とし，$0$ の開近傍から $M$ への滑らかな曲線
$\gamma:(-\varepsilon,\varepsilon)\to M$ で
$\gamma(0)=p$ となるものを考える．
二曲線 $\gamma_1,\gamma_2$ が同じ接ベクトルを表すとは，
$p$ を含む任意の座標写像 $\varphi$ について
\[
 (\varphi\circ\gamma_1)'(0)
 =(\varphi\circ\gamma_2)'(0)
\]
となることをいう．
その同値類を $T_pM$ と書き，$p$ における接空間という．
\end{definition}

\begin{lemma}[接ベクトルの座標表示]\label{lem:tangent-coord}
座標近傍 $(U,\varphi)$ と $p\in U$ を固定すると，
\[
 T_pM\longrightarrow\mathbb R^n,\qquad
 [\gamma]\longmapsto(\varphi\circ\gamma)'(0)
\]
は全単射である．
この同一視によって $T_pM$ は自然に
$n$ 次元実線形空間となる．
\end{lemma}
\begin{proof}
同値関係の定義から写像は単射である．
任意の $v\in\mathbb R^n$ に対し，
$t$ が十分小さいとき
$\varphi(p)+tv\in\varphi(U)$ となるので
\[
 \gamma_v(t)=\varphi^{-1}(\varphi(p)+tv)
\]
を定められる．
この曲線の座標速度は $v$ なので全射である．
$\mathbb R^n$ の線形構造を引き戻して
$T_pM$ の線形構造を定める．
別の座標 $\psi$ に取り替えると速度は
\[
 (\psi\circ\gamma)'(0)
 =D(\psi\circ\varphi^{-1})_{\varphi(p)}
  (\varphi\circ\gamma)'(0)
\]
と変換される．
Jacobi行列は線形同型なので，
定義した線形構造は座標の選択によらない．
\end{proof}

\section{導分としての接ベクトル}

\begin{definition}
$p\in M$ の近傍で定義された滑らかな実数値関数の芽全体を
$C_p^\infty(M)$ と書く．
実線形写像
$D:C_p^\infty(M)\to\mathbb R$ が
\[
 D(fg)=f(p)D(g)+g(p)D(f)
\]
を満たすとき，$p$ における導分という．
\end{definition}

\begin{lemma}[導分の局所表示]\label{lem:derivation}
座標近傍 $(U,\varphi)$，$\varphi=(x^1,\dots,x^n)$ と
$p\in U$ を固定する．
任意の導分 $D$ は一意に
\[
 D(f)=\sum_{j=1}^na_j\frac{\partial(f\circ\varphi^{-1})}
 {\partial x^j}(\varphi(p))
\]
と表される．
係数は $a_j=D(x^j)$ である．
\end{lemma}
\begin{proof}
まず $D(1)=D(1\cdot1)=2D(1)$ より $D(1)=0$ である．
従って定数関数への作用は零である．
Euclid空間の滑らかな関数 $g$ に対し，
$a=\varphi(p)$ として
\[
 g(x)-g(a)=\sum_{j=1}^n(x^j-a^j)h_j(x),
 \qquad
 h_j(x)=\int_0^1\partial_jg(a+t(x-a))\,dt
\]
と書ける．
$h_j$ は滑らかで $h_j(a)=\partial_jg(a)$ である．
この等式を $M$ の座標近傍へ引き戻し，
導分のLeibniz則を適用すると
\[
 D(f)=\sum_jD(x^j)h_j(a)
 =\sum_jD(x^j)\partial_jg(a).
\]
従って表示式を得る．
係数は座標関数 $x^j$ へ代入することで一意である．
\end{proof}

\begin{theorem}[曲線と導分による接空間の同一視]
\label{thm:tangent-derivation}
写像
\[
 T_pM\longrightarrow\operatorname{Der}(C_p^\infty(M),\mathbb R),
 \qquad
 [\gamma]\longmapsto
 \left(f\longmapsto\frac d{dt}f(\gamma(t))\bigg|_{t=0}\right)
\]
は線形同型である．
\end{theorem}
\begin{proof}
曲線の同値類が同じなら座標速度が同じであり，
連鎖律から全ての $f$ に対する導関数は同じである．
従って写像はwell-definedである．
曲線の導関数は和について線形であり，
積の微分公式からLeibniz則を満たすので導分となる．
座標近傍を選ぶと，
$[\gamma]$ は速度ベクトル $v$ に対応し，
写像の像は
\[
 D_v(f)=\sum_jv^j\partial_j(f\circ\varphi^{-1})(\varphi(p))
\]
となる．
補題\ref{lem:derivation}によって全ての導分は
この形で一意に表されるので，
写像は線形全単射である．
\end{proof}

\section{接写像と函手性}

\begin{definition}
滑らかな写像 $F:M\to N$ と $p\in M$ に対して
\[
 dF_p:T_pM\to T_{F(p)}N,\qquad
 dF_p([\gamma])=[F\circ\gamma]
\]
を接写像という．
\end{definition}

\begin{theorem}[接写像の連鎖律]\label{thm:chain}
滑らかな写像 $F:M\to N$，$G:N\to P$ について
\[
 d(G\circ F)_p=dG_{F(p)}\circ dF_p,\qquad
 d(\operatorname{id}_M)_p=\operatorname{id}_{T_pM}
\]
が成立する．
\end{theorem}
\begin{proof}
任意の $v=[\gamma]\in T_pM$ について
\[
 d(GF)_p(v)=[G\circ F\circ\gamma]
 =dG_{F(p)}([F\circ\gamma])
 =(dG_{F(p)}\circ dF_p)(v)
\]
となる．
恒等写像の場合は
$d(\operatorname{id}_M)_p([\gamma])=[\gamma]$
なので従う．
\end{proof}

\begin{proposition}[座標による接写像]\label{prop:jacobian}
座標写像 $\varphi$，$\psi$ のもとで
$dF_p$ は
\[
 D(\psi\circ F\circ\varphi^{-1})_{\varphi(p)}
\]
というJacobi行列によって表される．
\end{proposition}
\begin{proof}
$v=[\gamma]$ とする．
\[
 (\psi\circ F\circ\gamma)'(0)
 =D(\psi\circ F\circ\varphi^{-1})_{\varphi(p)}
  (\varphi\circ\gamma)'(0)
\]
がEuclid空間の連鎖律から成立する．
左辺は $dF_p(v)$ の $\psi$ 座標，
右辺の最後は $v$ の $\varphi$ 座標だから
求める表示式である．
\end{proof}

\section{埋め込み・沈め込みと階数定理}

\begin{definition}
$F:M^m\to N^n$ が各点で
$dF_p$ の階数 $m$ を持つとき\emph{はめ込み}，
階数 $n$ を持つとき\emph{沈め込み}という．
はめ込みであり，さらにその像への同相写像であるものを
\emph{埋め込み}という．
\end{definition}

\begin{theorem}[多様体の階数定理]\label{thm:rank}
滑らかな写像 $F:M^m\to N^n$ が $p$ の近傍で
一定の階数 $r$ を持つとする．
このとき $p$ と $F(p)$ の座標近傍を適切に選べば，
$F$ の座標表示は
\[
 (x_1,\dots,x_m)\longmapsto
 (x_1,\dots,x_r,0,\dots,0)
\]
となる．
\end{theorem}
\begin{proof}
$p$ を含む座標近傍と $F(p)$ を含む座標近傍を選び，
$F$ をEuclid空間の開集合間の滑らかな写像
$f:\mathbb R^m\supset U\to V\subset\mathbb R^n$ として表す．
座標変換は微分が可逆なので，
$f$ は対応する点の近傍で階数 $r$ を持つ．
Euclid空間の定階数定理によって，
両側の局所微分同相写像を合成することで
$f$ は上の標準形へ変換できる．
それらの局所微分同相写像を元の座標写像と合成すれば，
$M,N$ の新しい座標近傍となり所望の表示を得る．
\end{proof}

\begin{corollary}[正則値の逆像]\label{cor:regular-value}
$F:M^m\to N^n$ が滑らかで，
$q\in N$ が正則値，すなわち全ての $p\in F^{-1}(q)$ で
$dF_p$ が全射であるとする．
このとき $F^{-1}(q)$ は
$M$ の埋め込まれた $(m-n)$ 次元部分多様体であり，
\[
 T_p(F^{-1}(q))=\ker dF_p
\]
が成り立つ．
\end{corollary}
\begin{proof}
正則値の仮定から各 $p\in F^{-1}(q)$ で
$dF_p$ は階数 $n$ を持つ．
全射線形写像の近傍では適切な小行列式が非零であり続けるので，
階数は $p$ の近傍で $n$ のままである．
定理\ref{thm:rank}を適用すると，
局所座標において $F$ は
$(x_1,\dots,x_m)\mapsto(x_1,\dots,x_n)$ となる．
従って $F^{-1}(q)$ は
$x_1=\cdots=x_n=0$ によって局所的に記述されるため，
$\mathbb R^{m-n}$ と同相な埋め込み部分多様体となる．
その接空間は最初の $n$ 座標が零のベクトル全体であり，
これは座標表示の微分の核に等しい．
\end{proof}

\section{接束とベクトル場}

\begin{definition}
接空間の非交和
\[
 TM=\bigsqcup_{p\in M}T_pM
\]
を接束といい，
$\pi:TM\to M$，$\pi(v)=p$ を射影という．
滑らかな写像 $X:M\to TM$ で $\pi\circ X=\operatorname{id}_M$
を満たすものを滑らかなベクトル場という．
\end{definition}

\begin{proposition}[接束の局所自明性]\label{prop:tangent-bundle}
$M$ が $n$ 次元滑らかな多様体なら，
$TM$ は自然に $2n$ 次元滑らかな多様体となり，
各座標近傍 $(U,\varphi)$ に対して
\[
 \pi^{-1}(U)\longrightarrow\varphi(U)\times\mathbb R^n,
 \qquad
 v\in T_pM\longmapsto(\varphi(p),d\varphi_p(v))
\]
は微分同相写像となる．
\end{proposition}
\begin{proof}
表示した写像を $\Phi_\varphi$ と書く．
補題\ref{lem:tangent-coord}によって
各ファイバー上で全単射であるから，
$\Phi_\varphi$ は集合として全単射である．
これらの局所写像を座標図とする位相と
滑らかな構造を $TM$ に入れる．
二座標 $\varphi,\psi$ の遷移写像は
\[
 (x,v)\longmapsto
 \bigl(\psi\circ\varphi^{-1}(x),
       D(\psi\circ\varphi^{-1})_xv\bigr)
\]
であり，座標変換の滑らかさから滑らかである．
その逆は座標を交換したものなので滑らかである．
従ってこれらの座標は両立し，
$TM$ に $2n$ 次元の滑らかな構造が定まる．
射影 $\pi$ は局所座標で $(x,v)\mapsto x$ となり滑らかである．
\end{proof}

\section{単位の分割}

\begin{definition}
開被覆 $\{U_i\}_{i\in I}$ に従属する
\emph{滑らかな単位の分割}とは，
滑らかな関数族 $\{\rho_i:M\to[0,1]\}$ で
各 $\operatorname{supp}\rho_i\subset U_i$ を満たし，
族 $\{\operatorname{supp}\rho_i\}$ が局所有限であり，
\[
 \sum_i\rho_i(p)=1\qquad(p\in M)
\]
となるものをいう．
\end{definition}

\begin{theorem}[単位の分割の存在]\label{thm:partition}
第二可算Hausdorffな滑らかな多様体の任意の開被覆に対し，
その被覆に従属する滑らかな単位の分割が存在する．
\end{theorem}
\begin{proof}
第二可算Hausdorff多様体は局所コンパクトかつ
パラコンパクトであるという一般位相の定理を用いる．
従って与えられた開被覆に対し，
座標近傍の相対コンパクトな開集合からなる
局所有限な細分 $\{V_j\}$ を選び，
さらに
$\overline{W_j}\subset V_j$ となる開集合族
$\{W_j\}$ が全体を覆うようにできる．
この事実はパラコンパクト性による
縮小可能な局所有限開細分の存在定理から従う．
各 $j$ について座標写像 $\varphi_j$ を選び，
Euclid空間上の滑らかなコンパクト台関数
$\beta_j\ge0$ で
$\beta_j>0$ on $\overline{W_j}$，
$\operatorname{supp}\beta_j\subset V_j$ となるものを構成する．
これは $\overline{W_j}$ を有限個の小さい座標球で覆い，
各閉球上で正となる標準的な滑らかなbump関数を
有限和で組み合わせればよい．
座標外では零と延長すれば $M$ 上の滑らかな関数となる．
局所有限性から
\[
 B(p)=\sum_j\beta_j(p)
\]
は滑らかな関数であり，
$\{W_j\}$ が全体を覆い，各 $\beta_j$ が対応する
$\overline{W_j}$ 上で正なので $B(p)>0$ である．
従って
$\rho_j(p)=\beta_j(p)/B(p)$ と置けば滑らかで非負であり，
局所有限性と $\sum_j\rho_j=1$ を満たす．
各 $\rho_j$ の台は $V_j$ に含まれ，
細分により与えられた開被覆のいずれかの元に含まれる．
同じ被覆の元に対応する $\rho_j$ を必要に応じて
まとめれば，元の被覆に従属する分割が得られる．
\end{proof}

\begin{remark}
単位の分割は，Euclid空間の局所的な議論を
多様体上の大域的な議論へ貼り合わせる基本的な道具である．
ただしその存在にはパラコンパクト性が必要であり，
ここでは第二可算Hausdorff多様体の一般位相的性質を
明示的に前提とした．
\end{remark}

\section{到達課題の完全解答}
\noindent\textbf{球面の接空間を求める．}
\begin{proposition}
$S^2=\{p\in\mathbb R^3:|p|^2=1\}$ の
$p\in S^2$ における接空間は
\[
T_pS^2=\{v\in\mathbb R^3:p\cdot v=0\}
\]
である．
\end{proposition}
\begin{proof}
$F:\mathbb R^3\to\mathbb R$，
$F(x)=|x|^2$ と置くと
$dF_p(v)=2p\cdot v$ である．
$p\ne0$ だから $dF_p$ は全射であり，
1は正則値である．
正則値定理から $S^2=F^{-1}(1)$ は
二次元部分多様体で，
$T_pS^2=\ker dF_p=p^\perp$ となる．
\end{proof}
\noindent\textbf{到達確認．}
北極 $p=(0,0,1)$ の接空間を求めよ．
\par\noindent\textbf{解答．}
$p\cdot v=v_3$ なので
$T_pS^2=\{(v_1,v_2,0)\}$ である．
\end{document}